\section{Остановка и сходимость}
\label{sec:rates}
Шаг использует наблюдаемый зазор и не требует $\gamma$.
Известное значение $\gamma$ нужно для нижней границы в
\eqref{eq:stop}. Из квазар-выпуклости следует
\begin{equation}\label{eq:retained_bound}
 0\le f(x_k)-f(x^*)\le f(x_k)-\ell_k\le G(x_k)/\gamma.
\end{equation}

\begin{theorem}[Конечная остановка с гарантией точности]
\label{thm:gap_stopping}
При условиях \eqref{eq:gap_smoothness} и \eqref{eq:quasar}
алгоритм останавливается с $f(x_N)-f(x^*)\le\varepsilon$.
Если $G(x_0)>\gamma\varepsilon$, положим
$J=\lceil\log_2(G(x_0)/(\gamma\varepsilon))\rceil$.
Число шагов удовлетворяет
\begin{equation}\label{eq:gap_iteration_bound}
 N\le\left\lceil
 \frac{\max\{4,4L_1D^2/\log 2\}}{\gamma}\,J
 +\frac{16L_0D^2}{(\log 2)\gamma^2\varepsilon}
 \right\rceil.
\end{equation}
Вычислений функции, градиента и запросов линейной минимизации
ровно по $N+1$, включая последний запрос.
Если $G(x_0)\le\gamma\varepsilon$, алгоритм останавливается сразу.
\end{theorem}
\begin{proof}
При $L_0=L_1=0$ функция аффинна и нужен не более чем один полный шаг.
В остальных случаях из \eqref{eq:step_lower} и \eqref{eq:descent} следует
\begin{equation}\label{eq:gap_progress}
 f(x_k)-f(x_{k+1})\ge
 \min\left\{\frac{G(x_k)}2,
 \frac{(\log 2)G(x_k)^2}{2D^2(L_0+L_1G(x_k))}\right\}.
\end{equation}
Значение функции убывает, а $\ell_k$ возрастает, поэтому
$f(x_k)-\ell_k$ убывает. Положим $b_j=G(x_0)/(\gamma2^j)$,
$0\le j<J$. Каждый шаг до остановки отнесём к фазе
$b_j/2<f(x_k)-\ell_k\le b_j$.
В этой фазе \eqref{eq:retained_bound} даёт $G(x_k)>\gamma b_j/2$.
Функция $t^2/(L_0+L_1t)$ возрастает при $t>0$, поэтому каждый
такой шаг уменьшает значение функции не меньше чем на
\[
 \min\left\{\frac{\gamma b_j}{4},
 \frac{(\log 2)\gamma^2b_j^2}{8D^2(L_0+L_1\gamma b_j/2)}\right\}.
\]
Суммарное уменьшение на шагах, начинающихся в фазе, не больше
истинной ошибки при входе, то есть не больше $b_j$.
Сюда входит и последний шаг, выводящий из фазы: его конечная
точка допустима. Поэтому число шагов в фазе не больше
\[
 \max\left\{\frac4\gamma,
 \frac{8L_0D^2}{(\log 2)\gamma^2b_j}
 +\frac{4L_1D^2}{(\log 2)\gamma}\right\}.
\]
Суммирование и $\sum_{j<J}b_j^{-1}<2/\varepsilon$ доказывают
\eqref{eq:gap_iteration_bound}. Та же конечная граница верна для
каждого незавершённого начала процесса, поэтому бесконечный процесс
без остановки невозможен. В каждой посещённой точке по одному разу
вычисляются функция, градиент и линейный минимум.
Явный шаг \eqref{eq:log_step} не требует пробных вычислений функции.
\end{proof}

\subsection{Случай \texorpdfstring{$L_0=0$}{L0=0}}
\begin{corollary}[Геометрическая сходимость]\label{cor:pure_gap_linear}
Если $L_0=0<L_1$, положим
\[
 \underline\alpha_0=\min\{1,(\log 2)/(L_1D^2)\}.
\]
Для каждой вычисленной точки выполнено
\begin{align*}
 f(x_k)-f(x^*)&\le
 (1-\gamma\underline\alpha_0/2)^k\bigl(f(x_0)-f(x^*)\bigr),\\
 f(x_k)-\ell_k&\le
 (1-\gamma\underline\alpha_0/2)^k G(x_0)/\gamma.
\end{align*}
Если начальная проверка не дала остановки, достаточно не более
\[
 \left\lceil\frac{\log(G(x_0)/(\gamma\varepsilon))}
 {-\log(1-\gamma\underline\alpha_0/2)}\right\rceil
\]
шагов.
\end{corollary}
\begin{proof}
Граница шага даёт $\alpha_k\ge\underline\alpha_0$, поэтому
уменьшение функции не меньше $\underline\alpha_0G(x_k)/2$.
Для первой рекуррентной оценки используем \eqref{eq:quasar},
а для второй \eqref{eq:retained_bound} и $\ell_{k+1}\ge\ell_k$.
\end{proof}
Здесь $\nabla^2f(x^*)=0$ следует из $G(x^*)=0$.
Невыпуклый пример в разделе~\ref{sec:examples} обладает этим свойством
и ненулевым градиентом в минимуме.

\subsection{Сильно выпуклая допустимая область}
Используем определение из \cite{vyguzov}: множество $Q$ сильно
выпукло с константой $\lambda>0$, если для любых $u,v\in Q$
замкнутый шар радиуса $\lambda\|u-v\|^2$ с центром $(u+v)/2$
содержится в $Q$. Это условие на область, а не сильная выпуклость
целевой функции.

\begin{theorem}[Ненулевой градиент в минимуме и линейная скорость]
\label{thm:strong_set_linear}
Пусть дополнительно $Q$ сильно выпукло и $\nabla f(x^*)\ne0$.
Тогда минимум $x^*$ единственен и
\[
 c:=\min_{x\in Q}\|\nabla f(x)\|>0.
\]
Положим
\[
 \underline\alpha_{\rm sc}=
 \min\left\{1,\left(\frac{L_0}{2\lambda c}+2L_1D^2\right)^{-1}\right\},
\]
причём при нулевом знаменателе значение равно $1$. Тогда
\begin{align*}
 f(x_k)-f(x^*)&\le
 (1-\gamma\underline\alpha_{\rm sc}/2)^k
       \bigl(f(x_0)-f(x^*)\bigr),\\
 f(x_k)-\ell_k&\le
 (1-\gamma\underline\alpha_{\rm sc}/2)^k G(x_0)/\gamma.
\end{align*}
Алгоритм не использует ни $c$, ни $\lambda$; они нужны только
для этой оценки скорости.
\end{theorem}
\begin{proof}
Если бы $y$ было другим минимумом, условие \eqref{eq:quasar}
означало бы, что отношение
\[
 \frac{f(x^*+t(y-x^*))-f(x^*)}{t^\gamma}
\]
не убывает при $0<t\le1$. Оно неотрицательно и равно нулю при
$t=1$, поэтому функция постоянна на всём отрезке и
$\langle\nabla f(x^*),y-x^*\rangle=0$.
Однако условие на шар и необходимое условие минимума на $Q$ дают
\[
 \langle\nabla f(x^*),y-x^*\rangle
 \ge2\lambda\|\nabla f(x^*)\|\|y-x^*\|^2>0
\]
при $y\ne x^*$. Значит, минимум единственен.
Нулевой градиент в любой точке $Q$ по \eqref{eq:quasar} означал бы
ещё один минимум. Теперь непрерывность и компактность дают $c>0$.

Та же конструкция с шаром в точке линейного минимума даёт
\[
 G(x_k)\ge2\lambda\|\nabla f(x_k)\|\|d_k\|^2
           \ge2\lambda c\|d_k\|^2.
\]
При $L_0=L_1=0$ имеем $\alpha_k=1$. В остальных случаях используем
$\log(1+t)\ge t/(1+t)$ в \eqref{eq:log_step}
или указанный предел при $L_1=0$:
\[
 \alpha_k\ge\min\left\{1,
 \frac{G(x_k)}{(L_0+L_1G(x_k))\|d_k\|^2
                      +L_1D\|d_k\|G(x_k)}\right\}
 \ge\underline\alpha_{\rm sc}.
\]
Здесь $\|d_k\|\le D$, а предыдущее неравенство ограничивает
знаменатель, делённый на $G(x_k)$, числом
$L_0/(2\lambda c)+2L_1D^2$.
Обе рекуррентные оценки следуют так же, как в
следствии~\ref{cor:pure_gap_linear}.
\end{proof}
